\section{Experiment 1: Introduction to MySQL}
\label{exp:intro_mysql}

\subsection*{Aim}
To study the MySQL relational database management system, its architectural framework, core features, supported datatypes, and basic administrative commands.

\subsection*{Objectives}
\begin{itemize}[noitemsep,topsep=2pt]
    \item To understand the fundamentals and client-server architecture of MySQL RDBMS.
    \item To explore the classification of SQL statements (DDL, DML, DQL, DCL, TCL).
    \item To execute environment verification and administrative commands in the MySQL client.
\end{itemize}

\subsection*{Theory}
\textbf{MySQL} is an open-source Relational Database Management System (RDBMS) based on the client-server architecture. The server daemon (\texttt{mysqld}) manages storage engines, executes queries, and coordinates connections, while client programs submit queries and receive result sets.

\subsubsection*{1. Storage Engines and SQL Sublanguages}
MySQL features a pluggable storage engine architecture:
\begin{itemize}[noitemsep,topsep=2pt]
    \item \textbf{InnoDB:} Default ACID-compliant engine with row locking, crash recovery, and foreign keys.
    \item \textbf{MyISAM:} Fast engine optimized for read-intensive, non-transactional tasks.
\end{itemize}
SQL commands are broadly categorized into five functional sublanguages:
\begin{itemize}[noitemsep,topsep=2pt]
    \item \textbf{DDL (Data Definition Language):} Defines schemas (\texttt{CREATE}, \texttt{ALTER}, \texttt{DROP}).
    \item \textbf{DML (Data Manipulation Language):} Manages records (\texttt{INSERT}, \texttt{UPDATE}, \texttt{DELETE}).
    \item \textbf{DQL (Data Query Language):} Retrieves data tuples (\texttt{SELECT}).
    \item \textbf{DCL / TCL:} Manages permissions (\texttt{GRANT}, \texttt{REVOKE}) and transactions (\texttt{COMMIT}, \texttt{ROLLBACK}).
\end{itemize}

\subsubsection*{2. Core MySQL Data Types}
MySQL provides standard scalar data types to represent table attributes:
\begin{itemize}[noitemsep,topsep=2pt]
    \item \textbf{Numeric:} \texttt{INT} (integer), \texttt{DECIMAL(M,D)} (fixed-point), and \texttt{FLOAT} (floating-point).
    \item \textbf{String:} \texttt{CHAR(N)} (fixed-length), \texttt{VARCHAR(N)} (variable-length), and \texttt{TEXT}.
    \item \textbf{Date and Time:} \texttt{DATE} (\texttt{YYYY-MM-DD}), \texttt{TIME} (\texttt{HH:MM:SS}), and \texttt{TIMESTAMP}.
\end{itemize}

\newpage

\subsection*{Commands Execution \& Output Screenshots}

\subsubsection*{1. Connecting to MySQL Server and Verifying Session}
The MySQL CLI client is launched with root credentials, and session metadata is verified:

\begin{lstlisting}[caption={Authentication and version verification}]
-- Connect to MySQL server as root user
mysql -u root -p

-- Retrieve MySQL version, active user, and system timestamp
SELECT VERSION(), USER(), NOW();
\end{lstlisting}

\begin{figure}[H]
    \centering
    \includegraphics[width=0.88\linewidth,height=6.0cm,keepaspectratio]{Exp1-mysql-u.png}
    \caption{Connecting to MySQL monitor and verifying server version and session metadata.}
    \label{fig:exp1_mysql_u}
\end{figure}

\subsubsection*{2. Inspecting Server Runtime Status}
The \texttt{STATUS;} command queries active connection parameters and character sets:

\begin{lstlisting}[caption={Querying server status and connection parameters}]
-- Check active session connection details and server status
STATUS;
\end{lstlisting}

\begin{figure}[H]
    \centering
    \includegraphics[width=0.88\linewidth,height=6.0cm,keepaspectratio]{Exp1-status.png}
    \caption{Execution of \texttt{STATUS;} command displaying runtime parameters and character sets.}
    \label{fig:exp1_status}
\end{figure}

\newpage

\subsubsection*{3. Database Inspection, Creation, and Selection}
Existing catalogs are queried with \texttt{SHOW DATABASES;}. A sandbox database \texttt{lab\_demo} is created, selected with \texttt{USE}, and confirmed via \texttt{SELECT DATABASE()}:

\begin{lstlisting}[caption={Listing, creating, and switching active database}]
-- List all available databases on the server
SHOW DATABASES;

-- Create sandbox database and switch context
CREATE DATABASE lab_demo;
USE lab_demo;

-- Confirm active database context
SELECT DATABASE();
\end{lstlisting}

\begin{figure}[H]
    \centering
    \includegraphics[width=0.88\linewidth,height=9.0cm,keepaspectratio]{Exp1-use.png}
    \caption{Listing databases, creating \texttt{lab\_demo}, and verifying active catalog.}
    \label{fig:exp1_use}
\end{figure}

\subsection*{Result}
The installation and basic administrative commands of the MySQL RDBMS were studied and verified successfully. Client authentication, server status inspection, and database creation/selection were demonstrated.
